
在 KMP 中，我們預處理的方式是找出最長的重複前後綴，這樣在匹配失敗時可以跳過一些不必要的比較。

除了重複的前後綴以外，我們也可以使用 \textbf{最長公共前綴（LCP, Longest Common Prefix）} 的資訊來優化字串處理流程。
這樣的想法即是 \textbf{擴展KMP}（又稱 Z-Algorithm）背後的核心思想。    

\subsection{Z-陣列定義}

給定一個字串 $S$，我們定義 Z-陣列 $Z[i]$ 表示：

$$
Z[i] = \max \left\{ k \mid S[0 \ldots k-1] = S[i \ldots i+k-1] \right\}
$$

換句話說，$Z[i]$ 是從位置 $i$ 開始的字串與整個字串開頭（$S[0\ldots]$）的最長公共前綴長度。

舉例而言，設字串 $S=\texttt{ababc}$ 且 $i=2$，則 $Z[2] = 2$，因為從位置 $2$ 開始的字串 $\texttt{abc}$ 
與整個字串 $\texttt{ababc}$ 的最長公共前綴長度為 $2(\texttt{ab})$。

舉一個比較長的例子，對字串 $S = \texttt{aabcaabxaaaz}$，我們有：

\begin{center}
\begin{tabular}{c|c}
$i$ & $Z[i]$ \\
\hline
0 & -（不定義或取整體長度）\\
1 & 1 \\
2 & 0 \\
3 & 0 \\
4 & 3 \\
5 & 1 \\
6 & 0 \\
7 & 0 \\
8 & 2 \\
9 & 2 \\
10 & 1 \\
11 & 0 \\
\end{tabular}
\end{center}

\subsection{Z-陣列的構造}

我們可以在 $O(n)$ 的時間內構造整個 $Z$ 陣列，透過維護一個目前可延伸的區間 $[L, R]$，代表目前為止右端延伸最遠的 prefix 相等區段。

\lstset{caption={建構Z陣列的 C++ 實作}}
\begin{lstlisting}[language=C++]
vector<int> Z_Function(const string& S) {
    int n=S.size();
    vector<int> Z(n, 0);
    int L=0, R=0;

    for (int i=1;i<n;i++){
        if(i<=R)
            Z[i]=min(R-i+1, Z[i-L]);

        while(i+Z[i]<n && S[Z[i]]==S[i+Z[i]])
            Z[i]++;

        if(i+Z[i]-1>R){
            L=i;
            R=i+Z[i]-1;
        }
    }

    return Z;
}
\end{lstlisting}

當 $i > R$，我們只能從頭一個一個比對；但如果 $i \le R$，我們可以利用對稱性嘗試從 $Z[i - L]$ 的資訊推進。


\subsection{Z-陣列的應用}

Z 陣列的核心優勢是可以快速回答以下問題：

\begin{itemize}
    \item 一個 pattern 是否出現在 text 中
    \item 所有位置的最長公共前綴長度
    \item 判斷字串是否具有週期性（最短 border）
    \item 對於 $S = P + \texttt{\$} + T$，只要找到 $Z[i] = |P|$ 的位置，即可得知 pattern $P$ 在 text $T$ 中出現的位置
\end{itemize}

\subsection{與 KMP 的比較}

Z 陣列與 KMP prefix function 本質上都處理字串的重複結構，但：

\begin{itemize}
    \item KMP 的 prefix function 關心的是 $S[0 \ldots i]$ 中的最大前綴等於後綴的長度
    \item Z 陣列則關心每個 $i$ 開頭與整體字串最長前綴的匹配程度
    \item KMP 適合用在動態比對中（遇到 mismatch 立即跳轉）
    \item Z 陣列更適合整體前綴匹配分析與 pattern + text 一次合併處理
\end{itemize}

\subsection{總結}

Z 陣列提供了另一種觀察字串內部結構的方式，其應用面向廣泛，特別適合與 pattern matching、週期性判斷、最長公共前綴等問題結合使用。
